%% HXMT HE 1B 饱和分析与光变曲线重建
%% 初稿：中文
\documentclass[aps,prd,twocolumn,nofootinbib,superscriptaddress]{revtex4-2}

\usepackage{amsmath,amssymb}
\usepackage{graphicx}
\usepackage{booktabs}
\usepackage{multirow}
\usepackage{xcolor}
\usepackage{xeCJK}
\setCJKmainfont{STSong}

%% 图片占位宏
\newcommand{\figplaceholder}[1]{%
  \begin{center}
    \framebox[0.9\columnwidth]{\parbox{0.85\columnwidth}{\centering\vspace{1cm}#1\vspace{1cm}}}
  \end{center}
}

\begin{document}

\title{慧眼卫星高能望远镜 1B 级数据饱和分析与光变曲线重建}

\author{郭昊轩}
\affiliation{中国科学院高能物理研究所，北京 100049}

\author{待定}
\affiliation{中国科学院高能物理研究所，北京 100049}

\begin{abstract}
慧眼卫星（HXMT）高能 X 射线望远镜（HE）在极高计数率事件（如明亮伽马射线暴）期间，
前端电子学的 FIFO 缓冲区溢出会导致部分物理事例丢失，产生饱和效应。
现有的 1K 级数据管线对饱和区域的处理能力有限，而更底层的 1B 级原始数据保留了完整的
CCSDS 包结构和硬件时间戳信息，为饱和区域的精细分析提供了可能。
本文发展了一套完整的 1B 级数据处理方法：
（1）基于 wrap tracking 和 SEC 锚点的四遍时间重建算法，将 19-bit ptime 计数器
恢复为绝对任务消逝时间（MET），在非饱和区与 1K 管线结果完美一致；
（2）饱和区间检测算法，涵盖 FIFO 复位整包丢失、包内无痕丢数和深度饱和三种模式；
（3）基于三机箱交叉参考的光变曲线重建方法，利用未饱和机箱的事件分布指导饱和区域补全。
利用三个代表性伽马射线暴（GRB~200415A、GRB~260226A、GRB~221009A）验证表明：
非饱和区 1B/1K 事件数差异 $\Delta = 0$；
极端饱和情况下（GRB~221009A），1B 重建保留了更完整的事件集（97.9\%--100\% 精确匹配率），
三机箱交叉参考覆盖率达 92\%。
该方法为利用 HXMT HE 1B 数据开展饱和区域的科学分析提供了可靠工具。
\end{abstract}

\maketitle

%% ======================================================================
\section{引言}\label{sec:intro}

慧眼卫星（Insight-HXMT）是中国首颗 X 射线天文卫星，搭载了高能 X 射线望远镜
（HE，20--250~keV）、中能 X 射线望远镜（ME，5--30~keV）和低能 X 射线望远镜
（LE，1--15~keV）三台载荷。
HE 望远镜拥有 5100~cm$^2$ 的有效面积，是目前在轨运行的最大面积硬 X 射线探测器之一，
在伽马射线暴（GRB）等高计数率瞬变源的观测中具有独特优势。

然而，HE 望远镜的前端电子学采用 FIFO 缓冲区架构进行数据中转，当探测器计数率超过
MCU 处理能力时，FIFO 缓冲区溢出会导致物理事例数据丢失。这种饱和效应在明亮 GRB
的观测中尤为突出，例如 GRB~221009A（史上最亮 GRB）期间，大量事例因 FIFO 溢出而丢失，
严重影响了光变曲线的完整性和科学分析的可靠性。

目前 HXMT 的标准数据处理管线（1K 级）在饱和区域的处理能力有限：
1K 管线对 FIFO 复位边界的事件做了较激进的过滤，可能额外丢失有效数据。
更底层的 1B 级原始数据保留了完整的 CCSDS 包结构、硬件时间戳和所有通过 CRC 校验的事件，
为饱和效应的精细分析和光变曲线重建提供了更丰富的信息。

本文发展了一套完整的 HXMT HE 1B 级数据饱和分析与光变曲线重建方法，主要贡献包括：
\begin{enumerate}
  \item 四遍时间重建算法，在 FIFO 拥塞和复位条件下准确恢复事件时间；
  \item 三种饱和模式的检测方法（FIFO 复位、包内无痕丢数、深度饱和）；
  \item 基于三机箱交叉参考的光变曲线重建策略；
  \item 利用三个不同饱和程度的 GRB 进行系统验证。
\end{enumerate}

本文结构如下：第~\ref{sec:instrument}~节介绍 HE 硬件架构与数据通路；
第~\ref{sec:satmech}~节分析饱和丢数的物理机制；
第~\ref{sec:method}~节详述算法方法；
第~\ref{sec:results}~节给出验证结果；
第~\ref{sec:discussion}~节讨论局限性与适用范围；
第~\ref{sec:conclusion}~节总结全文。

%% ======================================================================
\section{仪器与数据}\label{sec:instrument}

\subsection{HE 望远镜硬件架构}

HE 望远镜包含 18 个 NaI(Tl)/CsI(Na) 复合探测器，分置于三个独立机箱
（Box~A、Box~B、Box~C），每个机箱有 6 个探测器模块。
三个机箱拥有完全独立的数据通路，包括独立的 FPGA、FIFO 缓冲区和 MCU 控制器。

%% Figure 1: 数据通路架构图
\begin{figure}[t]
\figplaceholder{图~1：HE 数据通路架构图。
NaI/CsI 探测器 $\rightarrow$ ASIC $\rightarrow$ FPGA $\rightarrow$ FIFO~A (M67204H)
$\rightarrow$ MCU (8051) $\rightarrow$ FIFO~B $\rightarrow$ 1553B 总线 $\rightarrow$ 下传。
三个机箱（A/B/C）拥有完全独立的数据通路。}
\caption{HE 望远镜数据通路架构。物理事例从探测器经 ASIC 数字化后由 FPGA 写入
硬件 FIFO~A 缓冲区，MCU 从 FIFO~A 逐事件读取并打包为 CCSDS 包写入 FIFO~B，
最终经 1553B 总线下传。}
\label{fig:datapath}
\end{figure}

\subsection{数据通路}

探测器接收 X 射线/$\gamma$ 射线光子后产生模拟信号，经前端 ASIC 电路数字化为
事例数据（包含能道、精确时间等信息），由 FPGA 写入硬件 FIFO~A 缓冲区（M67204H 芯片）。

FIFO~A 是环形缓冲区结构，容量约 4096 字节（可容纳约 455 个事例）。
FIFO~A 满时不覆盖旧数据——FPGA 端的写入被阻塞，新事例直接丢弃。
FIFO~A 具有硬件 Full 标志位（\texttt{FIFOAFull1}、\texttt{FIFOAFull2}），
MCU 可通过写入 \texttt{ResetFIFOA = 0x00} 复位整个 FIFO。

MCU 是 8051 架构的单线程处理器，通过主循环轮询处理所有任务。
物理数据处理函数 \texttt{HandlePhysicalLVDS()} 每次从 FIFO~A 读取 109 个事例，
打包为一个 CCSDS 包写入 FIFO~B，耗时约 7~ms。
每个 CCSDS 包的结构为：6 字节包头 + 872 字节载荷（$109 \times 8$ 字节事例）
+ 4 字节 UTC 时间尾 = 882 字节。

MCU 主循环还包含 FIFO~A 满溢检查函数 \texttt{FIFOAFullReset()}：
当检测到 FIFO~A 的 Full 标志有效时，执行复位清空整个 FIFO~A 缓冲区。

\subsection{事例格式与时间戳}

每个物理事例占 8 字节，通过 4-bit CRC 校验和类型标志位区分三种类型：

\begin{itemize}
  \item \textbf{Event（物理事例）}：包含能道 channel 和 19-bit 精确时间计数器 ptime，
    分辨率 2~$\mu$s/tick，回绕周期 $2^{19} \times 2~\mu\text{s} = 1.048576$~s；
  \item \textbf{Second（秒标记）}：硬件每秒插入，记录整秒 MET（stime）和对应 ptime，
    构成时间重建的绝对锚点；
  \item \textbf{Error}：CRC 校验失败的事例。
\end{itemize}

MCU 读取事例时通过搜索 0x5A 同步字节定位事例边界：
检查上次读取的末尾字节是否为 0x5A，若非则逐字节扫描 FIFO~A 直到找到
（最多 2048 次），随后读取 8 字节事例数据。

\subsection{1B 与 1K 数据级别}

HXMT 的数据处理分为多个级别。\textbf{1B 级}是最底层的原始科学数据，
保留了完整的 CCSDS 包结构和所有硬件信息。
\textbf{1K 级}是经过时间校正、事件筛选等处理后的科学数据，
是大多数科学分析使用的标准级别。

1K 管线对事件进行了时间排序和部分过滤，在 FIFO 复位边界处采用了较激进的
事件剔除策略。本文的 1B 级处理则保留所有通过 CRC 校验的事件，
在时间重建后再进行饱和判断和光变曲线重建。

%% ======================================================================
\section{饱和丢数机制}\label{sec:satmech}

HE 望远镜的饱和丢数源于 FIFO~A 缓冲区的溢出。
当探测器计数率超过 MCU 的处理速度时，FIFO~A 中的数据积压最终导致缓冲区满溢，
新事例被丢弃。根据溢出持续时间和 MCU 主循环状态的不同，
产生两种截然不同的丢数模式。

%% Figure 2: 两种丢数模式时序示意图
\begin{figure}[t]
\figplaceholder{图~2：两种丢数模式时序示意图。
上方：模式一（包内无痕丢数），FIFO 短暂满溢后恢复，MCU 回到主循环时已不满。
下方：模式二（FIFO 复位整包丢失），FIFO 持续满载，MCU 回到主循环时仍满，
触发 FIFOAFullReset() 清空全部缓冲数据。}
\caption{HE 望远镜两种饱和丢数模式的时序对比。}
\label{fig:two_modes}
\end{figure}

\subsection{模式一：包内无痕丢数}

当物理事例率略高于 MCU 读取速率时，MCU 在执行 \texttt{HandlePhysicalLVDS()}
期间（约 7~ms），FPGA 的写入速度超过 MCU 的读出速度，FIFO~A 在某一时刻被写满。
此时 FPGA 无法写入，新到达的事例被直接丢弃。
MCU 继续从 FIFO~A 读取数据，腾出空间后 FPGA 恢复写入。

\texttt{HandlePhysicalLVDS()} 完成后 MCU 回到主循环，
此时 FIFO~A 已不再满，\texttt{FIFOAFullReset()} 检查 Full 标志位时
发现 FIFO 未满，不触发复位。

这种模式的关键特征是：
\begin{itemize}
  \item CCSDS 包结构完整，包间时间间隔无异常；
  \item 丢失少量事例（数十到数百个），在包内产生异常大的时间间隔；
  \item \textbf{无法通过包间 gap 检测}，但可通过包内事件间隔的统计异常发现。
\end{itemize}

\subsection{模式二：FIFO 复位整包丢失}

当物理事例率远高于 MCU 读取速率时，在 \texttt{HandlePhysicalLVDS()} 的 7~ms
执行期间，FPGA 的写入量远超 MCU 的读出量。
MCU 处理完一个 CCSDS 包回到主循环时，FIFO~A \textbf{仍然是满的}。
\texttt{FIFOAFullReset()} 检测到 Full 标志有效后执行
\texttt{ResetFIFOA = 0x00}，清空整个 FIFO~A 缓冲区。

这种模式的特征是：
\begin{itemize}
  \item 丢失整个 FIFO 缓冲区的数据（约 455 个事例及后续积压事例）；
  \item 相邻 CCSDS 包之间出现远大于正常间隔的时间空洞；
  \item \textbf{可通过包间 gap 检测}。
\end{itemize}

%% Table 2: 两种模式对比
\begin{table}[t]
\centering
\caption{三种饱和模式的检测与重建方法}
\label{tab:modes}
\begin{tabular}{lccc}
\toprule
 & 包内无痕丢数 & FIFO 复位 & 深度饱和 \\
\midrule
触发条件 & $R \gtrsim R_\mathrm{MCU}$ & $R \gg R_\mathrm{MCU}$ & 持续 $R \gg R_\mathrm{MCU}$ \\
丢失量级 & 数十--数百 & $\sim$FIFO 容量 & 包内 gap \\
可检测性 & 泊松检验 & 包间 gap & 邻居率判定 \\
重建策略 & 交叉参考 & 交叉参考 & burst rate \\
\bottomrule
\end{tabular}
\end{table}

\subsection{深度饱和}

当探测器计数率持续远超 MCU 处理能力时，所有包都处于 FIFO 拥塞状态。
MCU 的读取模式表现为周期性的突发（burst）：每约 8~ms 读取 109 个事件，
实际读取过程在约 0.5~ms 内完成，其余约 7.5~ms 为主循环开销期间。
在此期间 FIFO 持续积累事件，可能反复满溢丢弃。

深度饱和区的特征是：所有邻居包的表观事件率都低于 MCU 读速率
（$\sim$15000~evt/s），因为包的时间跨度反映的是 MCU 主循环周期
而非物理事件率。这使得常规的泊松检验和邻居率参考方法失效，
需要专门的包级修正策略。

%% ======================================================================
\section{方法}\label{sec:method}

\subsection{时间重建算法}\label{sec:time_recon}

1B 级数据中每个物理事例仅有 19-bit 的 ptime 计数器
（回绕周期 $P_\mathrm{wrap} = 2^{19} \times 2~\mu\text{s} \approx 1.0486$~s），
需要确定相对于 SEC 锚点的回绕次数 $n_\mathrm{wraps}$ 才能恢复绝对 MET 时间。

时间重建公式为：
\begin{equation}\label{eq:met}
  t_\mathrm{MET} = t_\mathrm{anchor} + (n_\mathrm{wraps} \times P_\mathrm{MOD}
    + p - p_\mathrm{anchor}) \times \Delta t + t_\mathrm{corr}
\end{equation}
其中 $t_\mathrm{anchor}$ 和 $p_\mathrm{anchor}$ 来自最近的有效 Second 事件，
$P_\mathrm{MOD} = 2^{19} = 524288$，$\Delta t = 2~\mu$s，
$t_\mathrm{corr} = 4.0$~s 为 1B$\rightarrow$1K 经验时间校正。

在正常观测条件下，$n_\mathrm{wraps}$ 的确定是直接的。
但在 GRB 饱和期间，三个因素使问题复杂化：
（1）FIFO 拥塞导致事件滞留，UTC 尾时间戳严重滞后于事件真实时间；
（2）FIFO 复位后，复位前最后接受的 SEC 仍是 MCU 当前锚点，
在下一个整秒边界 FPGA 注入新 SEC 之前，事件都使用该过期锚点，
路径 2 的 $n_\mathrm{base}$ 取整可能差 $\pm 1$；
（3）过期锚路径估算 $n_\mathrm{base}$ 时，包的 median ptime 经过 anchor ptime 附近
会引起取整翻转，产生 $\pm 1$ 的 $n_\mathrm{base}$ 偏差。

为此，我们发展了四遍（4-pass）处理管线（图~\ref{fig:pipeline}）。

%% Figure 3: 4-pass 流程图
\begin{figure}[t]
\figplaceholder{图~3：时间重建 4-pass 处理管线流程图。
Pass~1: 逐包重建（三路径 + FIFO 复位后向修正）$\rightarrow$
Pass~2: FIFO 复位后时间倒退修正 $\rightarrow$
Pass~3: ptime 回绕边界 dip 修正 $\rightarrow$
Pass~4: 分段排序。}
\caption{1B 级数据时间重建的四遍处理管线。}
\label{fig:pipeline}
\end{figure}

\subsubsection{Pass 1：逐包时间重建}

对每个 CCSDS 包内的事件，按优先级选择三条计算路径之一：

\textbf{路径 1——Wrap Tracking（FIFO 拥塞专用）}：
当 SEC 锚点变"陈旧"（距最近 SEC 超过 35 个包）且处于 FIFO 拥塞期间时，
用包内 ptime 中位数追踪回绕次数。
FIFO 拥塞期间 MCU 从 FIFO 顺序读出事件，连续包的中位数 ptime 单调递增。
当中位数发生大负跳变（回绕），累加 \texttt{congestion\_wrap\_count}。
为修正锚点 ptime 与中位数 ptime 的相位差，引入相位校正：
\begin{equation}
  n_\mathrm{corr} = n_\mathrm{cong} + \delta_\mathrm{phase}
\end{equation}
其中 $\delta_\mathrm{phase} \in \{-1, 0, +1\}$，取决于 anchor 中位数与
anchor ptime 之差是否超过半周期。

包内各事件相对中位数的偏移 $\delta = p - p_\mathrm{median}$ 用于判断包内回绕：
$\delta < -P_\mathrm{MOD}/2$ 时加 1 次回绕，$\delta > P_\mathrm{MOD}/2$ 时减 1 次。

\textbf{路径 2——包级过期锚路径}：
FIFO 复位后锚点过期多个 wrap 周期，但事件是新鲜的（无 FIFO 延迟）。
利用 UTC 尾时间戳和中位数 ptime 估算基础回绕数 $n_\mathrm{base}$：
\begin{equation}
  n_\mathrm{est} = \mathrm{round}\!\left(
    \frac{t_\mathrm{utc} - t_\mathrm{anchor} - (p_\mathrm{median} - p_\mathrm{anchor}) \Delta t}
    {P_\mathrm{wrap}}
  \right)
\end{equation}
从 $\{n_\mathrm{est}-1, n_\mathrm{est}, n_\mathrm{est}+1\}$ 中选择使
$|t_\mathrm{MET} - t_\mathrm{utc}|$ 最小的值。

\textbf{路径 3——正常路径（阈值法）}：
锚点新鲜时（距 SEC $\leq$35 包），至多 1 次 wrap。
当 $p - p_\mathrm{anchor} < -\theta$ 时判定为正向回绕（事件 ptime 跨过 0 回到小值），
$p - p_\mathrm{anchor} > P_\mathrm{MOD} - \theta$ 时判定锚点 ptime 已跨过回绕边界而事件尚未
（实际差值应为 $p - p_\mathrm{anchor} - P_\mathrm{MOD}$，即小负值），
阈值 $\theta = 10000$ ticks（$\approx$20~ms）。
注意 ptime 由硬件计数器产生，在 MCU 顺序读取 FIFO 时严格单调递增，
不存在物理上的反向回绕。两种情况都是正向回绕在不同 anchor 位置下的表现。

\subsubsection{Pass 1 补充：FIFO 复位后向 SEC 修正}

FIFO 复位清空了 FIFO 中的所有数据，复位后没有旧 SEC 混入。
真正的问题是：复位前最后接受的 SEC 仍然是 MCU 的当前锚点。
直到下一个整秒边界 FPGA 注入新 SEC，这中间的事件都使用该过期锚点。
旧锚点的 $(t_\mathrm{anchor}, p_\mathrm{anchor})$ 本身是正确的，但距现在已过数秒
（多个 wrap 周期），路径 2 用 $t_\mathrm{utc}$ + round() 估算 $n_\mathrm{base}$
时取整可能差 $\pm 1$，导致事件时间系统性偏移约 1 个 $P_\mathrm{wrap} \approx 1.05$~s。

解决方案是等待后续出现新鲜 SEC 锚点后，反向修正之前缓冲的事件。
修正分三个阶段：

\textbf{Phase~1（置信度分类）}：
用新鲜 SEC 为每个缓冲事件计算后向 wrap 数和置信度 margin
（0 = 完全模糊，0.5 = 完全确定）。margin $> 0.05$ 的事件直接采用后向计算。

\textbf{Phase~2（前后对照）}：
对模糊事件，比较过期锚点的前向估计和新鲜锚点的后向估计，
选择与前向结果更一致的后向候选。
原理是：两个锚点有不同的 anchor ptime，因此 wrap 边界位置不同，
当一个方向模糊时另一个方向通常能给出正确方向的估计。

\textbf{Phase~3（应用修正）}：将修正后的 wrap 数统一应用。

\subsubsection{Pass 2：FIFO 复位后时间倒退修正}

MCU 写入顺序本来就是时间顺序（FIFO 严格先进先出）。
Pass 2 修正的不是文件乱序，而是过期锚路径的 $n_\mathrm{base}$ 取整偏差
使一批包被放到高一个 $P_\mathrm{wrap}$ 的位置。

算法扫描相邻"干净包"（span $< 0.3$~s）间的反向跳跃 $\approx -P_\mathrm{wrap}$，
反向查找属于高层级的包批次。通过 gap 判据区分真假倒退：
批次前有大 gap（$> 0.3$~s）为 FIFO 复位后的真错位，整批下移 $-P_\mathrm{wrap}$；
批次前无 gap（平滑衔接）为纯文件重排的假倒退，跳过不处理。

\subsubsection{Pass 3：ptime 回绕边界修正}

Pass 3 修正的凹坑源于过期锚路径的 $n_\mathrm{base}$ 取整不稳定。
ptime 严格递增意味着对新鲜锚点，$p - p_\mathrm{anchor} < 0$ 一定是 wrap（无歧义）。
凹坑的真正原因是：过期锚路径估算 $n_\mathrm{base}$ 时用 round() 取整，
当包的 median ptime 经过 anchor ptime 附近时，
$(p_\mathrm{median} - p_\mathrm{anchor})$ 跳变导致分子变化 $\sim$1.05~s，
取整翻转 $\pm 1$，使一批包的时间比前后邻居低一个 $P_\mathrm{wrap}$。

修正分两步：
（1）检测正向跳跃 $\approx +P_\mathrm{wrap}$，查找被高层级包包夹的低层级批次，
整批上移；
（2）对包内同时包含两个层级事件的"混合包"（span $\approx P_\mathrm{wrap}$），
按邻居包层级判定多数/少数簇，将少数簇对齐。

Pass 2 和 Pass 3 是对称互补的操作：Pass 2 修偏高批次（下移 $-P_\mathrm{wrap}$），
Pass 3 修偏低批次（上移 $+P_\mathrm{wrap}$）。
Pass 2 需要 gap 判据区分真错位和文件重排产生的假倒退，
Pass 3 不需要——凹坑不会是文件重排。

\subsubsection{Pass 4：分段排序}

Pass 4 排的不是文件中的原始包乱序，而是后向 SEC 修正改变 MET 后
在修正区域和相邻 wrap tracking 区域交界处产生的局部时间重叠。
对非后向修正区域的连续段分别排序，
后向修正区域保持不动，防止排序破坏精确修正的结果。

\subsection{饱和区间检测}\label{sec:detection}

\subsubsection{FIFO 复位检测}

对重建后的包时间序列，计算相邻包间的时间间隔 gap。
采用自适应阈值：
\begin{equation}\label{eq:gap}
  \mathrm{gap} > T_\mathrm{baseline} \times F_\mathrm{gap}
\end{equation}
其中 $T_\mathrm{baseline}$ 取前后包中平均事件间隔的较小值
（即事件率较高的包），$F_\mathrm{gap} = 100$。
同时要求前后包的最大事件率超过 MCU 读速率下限
$R_\mathrm{MCU} = 15000$~evt/s（略低于理论值 $109/6.9~\text{ms} \approx 15797$~evt/s），
否则 FIFO 不可能溢出。

对于深度饱和区域，相邻包可能都是拥塞包导致事件率被低估。
此时扩展到前后各 5 个包的窗口取最大事件率判断。

\subsubsection{静默丢数检测（包内无痕丢数）}

对每个高事件率包，利用包内事件间隔的泊松统计异常检测无痕丢数。

对于正常饱和区（邻居包事件率 $\geq R_\mathrm{MCU}$），
用包内过滤间隔（$< 1$~ms）估算局部事件率 $\lambda$。
对于拥塞宽包（span $> 3 \times$ 邻居中位 span），
用邻居包事件率估算 $\lambda$。

对每个事件间隔 $\Delta t$ 计算泊松概率：
\begin{equation}\label{eq:poisson}
  \log_{10} p = -\frac{\lambda \cdot \Delta t}{\ln 10}
\end{equation}
若 $\log_{10} p < -10$，标记为静默丢数候选。
丢失事件数估算为 $N_\mathrm{lost} = \lambda \cdot \Delta t - 1$。

\subsubsection{深度饱和识别}

当所有邻居包的事件率均低于 $R_\mathrm{MCU}$（即都是拥塞包）时，
判定为深度饱和区。此区域不进行逐间隔的泊松检测
（因为 $\lambda$ 无法从数据可靠估算），而是采用包级修正策略。

\subsection{光变曲线重建}\label{sec:reconstruction}

根据三种饱和模式的不同特征，分别采用对应的重建策略，
三种重建独立进行且都仅使用原始观测事件做参考。

\subsubsection{FIFO 复位 gap 补全}

每个 FIFO reset gap 对应恰好一次 FIFO 复位（复位后 FIFO 清空，
即使 $R_\mathrm{true} = 50000$~evt/s 也需约 9~ms 才能重新填满，
而主循环一轮仅需微秒级）。
丢失事件数：
\begin{equation}\label{eq:nlost}
  N_\mathrm{lost} = R_\mathrm{true} \times T_\mathrm{gap}
\end{equation}

$R_\mathrm{true}$ 的估算利用物理约束：FIFO 持续满载时，
包内事件间距反映 MCU 读速率而非真实率；
但 FIFO 复位后的 post-reset 包不存在满载阻塞，
其 span 直接反映 $R_\mathrm{true}$：
\begin{equation}\label{eq:rtrue}
  R_\mathrm{true} \approx \frac{109}{T_\mathrm{span,post\text{-}reset}}
\end{equation}
深度饱和区 post-reset 包本身也是拥塞包时，
改用邻近包的 burst rate（包内过滤间隔 $< 1$~ms 部分的事件率）。

事件时间分布通过三机箱交叉参考确定：
利用其他未饱和机箱的事件在 gap 时间范围内的分布构建 1~ms 分辨率的
形状函数（shape bins），跳过参考机箱自身的不可信区间
（FIFO reset gap、拥塞宽包、含异常间隔的包）。

多个参考机箱的形状取均值，按 gap 两端附近的事件数比值进行标定（calibration）。
形状函数归一化到 $N_\mathrm{lost}$ 后按比例分配事件。
当参考覆盖率 $< 30\%$ 时退化为均匀分配。

\subsubsection{静默丢数补全}

每个检测到的异常间隔，用 $N_\mathrm{lost} = \lambda \cdot \Delta t - 1$
估算丢失数，同样通过交叉参考构建时间分布。
参考覆盖率 $< 30\%$ 时均匀分配。

\subsubsection{深度饱和包级修正}

对深度饱和区的每个包，用包内 burst 部分
（连续间隔 $< 1$~ms 的紧密事件簇）的计数率作为 $R_\mathrm{true}$。
包内超过 1~ms 的大间隔视为 MCU 主循环开销期间的数据丢失：
\begin{equation}\label{eq:deep}
  N_\mathrm{fill} = R_\mathrm{true} \times \sum_{i: \Delta t_i > 1\,\text{ms}} \Delta t_i
\end{equation}
补全事件均匀分配在各个 gap 间隔中。

\subsection{三机箱交叉参考}\label{sec:crossref}

三个机箱拥有完全独立的 FIFO 和 MCU，但观测同一 X 射线源，
计数率成比例。当一个机箱饱和时，其他机箱可能仍有完好数据。

\subsubsection{不可信区间检测}

交叉参考的质量取决于参考机箱的数据可靠性。
以下时间区间被标记为不可信，参考时跳过：

\begin{enumerate}
  \item \textbf{FIFO reset gap}：整包丢失的时间段，无事件；
  \item \textbf{拥塞宽包}：包跨时 $> 3 \times$ 邻居中位跨时，
    事件分布反映 MCU 读取模式而非真实光变；
  \item \textbf{含异常间隔的包}：包内有泊松异常间隔
    （用邻居率做 $\lambda$ 判定），整个包不可信。
\end{enumerate}

\subsubsection{覆盖率统计}

以 GRB~221009A 为例，三机箱交叉参考的覆盖情况见表~\ref{tab:coverage}。
约 92\% 的 FIFO reset gap 至少有一个参考机箱可用，
仅约 8\% 的 gap 三个机箱同时饱和。

%% Table: 交叉参考覆盖率
\begin{table}[t]
\centering
\caption{GRB~221009A 三机箱交叉参考覆盖统计}
\label{tab:coverage}
\begin{tabular}{lccc}
\toprule
覆盖情况 & Box A & Box B & Box C \\
\midrule
两个参考 box 都未饱和 & 49.2\% & 38.3\% & 46.8\% \\
至少一个参考 box 未饱和 & 92.0\% & 92.0\% & 92.6\% \\
三个 box 同时饱和 & 8.0\% & 8.0\% & 7.4\% \\
\bottomrule
\end{tabular}
\end{table}

%% ======================================================================
\section{验证与结果}\label{sec:results}

我们选择三个具有不同饱和程度的 GRB 进行系统验证，
其基本参数和验证统计见表~\ref{tab:grbs}。

%% Table 1: 三个 GRB 的参数
\begin{table*}[t]
\centering
\caption{三个验证 GRB 的基本参数与 1B/1K 对比统计}
\label{tab:grbs}
\begin{tabular}{lcccccc}
\toprule
GRB & 饱和程度 & Box & 1B 事件数 & 1K 事件数 & $\Delta$ & $\Delta/N_\mathrm{1K}$ \\
\midrule
\multirow{3}{*}{200415A} & \multirow{3}{*}{轻度} & A & --- & --- & 0 & 0.0\% \\
 & & B & --- & --- & 0 & 0.0\% \\
 & & C & --- & --- & 0 & 0.0\% \\
\midrule
\multirow{3}{*}{260226A} & \multirow{3}{*}{中度} & A & --- & --- & $\sim$0 & 0.0\% \\
 & & B & --- & --- & $\sim$0 & 0.0\% \\
 & & C & --- & --- & $\sim$0 & 0.0\% \\
\midrule
\multirow{3}{*}{221009A} & \multirow{3}{*}{极端} & A & 3,857,418 & 3,857,310 & +108 & +0.003\% \\
 & & B & 3,800,510 & 3,797,574 & +2,936 & +0.08\% \\
 & & C & 3,745,229 & 3,619,371 & +125,858 & +3.5\% \\
\bottomrule
\end{tabular}
\end{table*}

\subsection{GRB~200415A（轻度饱和）}\label{sec:grb200415}

GRB~200415A 仅在峰值处存在很小的饱和区间，整体饱和程度轻微，
适合作为验证时间重建算法正确性的基准样本。

所有三个机箱的 1B 重建结果与 1K 管线完全一致：
事件数差异 $\Delta = 0$，所有 1~s 时间 bin 内的计数率完美匹配。
这证明了四遍时间重建算法即使在存在轻度饱和的情况下也能准确恢复事件时间。

%% Figure 4: 1B/1K 对比
\begin{figure}[t]
\figplaceholder{图~4：GRB~200415A 的 1B/1K 光变曲线对比。
三个机箱的 1B 重建时间与 1K 管线时间完美对齐（$\Delta = 0$）。}
\caption{轻度饱和区 1B/1K 验证。GRB~200415A（峰值处有轻微饱和）的光变曲线对比，
三个机箱结果完全一致。}
\label{fig:validation}
\end{figure}

\subsection{GRB~260226A（中度饱和）}\label{sec:grb260226}

GRB~260226A 是一个中度饱和事件（$T_0 \approx \mathrm{MET}~446726273$），
具有明显的 FIFO 复位 gap，但无静默丢数。

1B/1K 事件数差异仅约 2 个事件（+0.0\%），所有 bin 误差 0.0\%。
Pass~3（ptime 回绕边界修正）在三个机箱中分别触发了 7、14、13 个包的修正。

%% Figure 5: 饱和区间检测示例
\begin{figure}[t]
\figplaceholder{图~5：GRB~260226A 饱和区间检测结果。
上方：重建光变曲线，箭头标注 FIFO reset gap 位置。
下方：包间间隔时序图，红色虚线为检测阈值。}
\caption{中度饱和区间的 FIFO reset gap 检测结果示例。}
\label{fig:detection}
\end{figure}

\subsection{GRB~221009A（极端饱和）}\label{sec:grb221009}

GRB~221009A（2022 年 10 月 9 日）是有记录以来最明亮的伽马射线暴，
HE 探测器经历了严重且持续的饱和。

\subsubsection{整体统计}

三个机箱的 1B/1K 对比见表~\ref{tab:grbs}：
\begin{itemize}
  \item Box~A：1B 比 1K 多 108 个事件（+0.003\%）；
  \item Box~B：1B 比 1K 多 2,936 个事件（+0.08\%）；
  \item Box~C：1B 比 1K 多 125,858 个事件（+3.5\%），其中大部分差异
    集中在 $T+510$~s 附近，系 1K 管线数据丢失（非 1B 问题）。
\end{itemize}

\subsubsection{逐事件精确匹配分析}

在 FIFO 复位恢复区域对 1B 和 1K 进行逐 channel、逐事件的匹配
（时间容差 $< 0.5$~ms），结果显示：
\begin{itemize}
  \item 精确匹配率：97.9\%--100\%；
  \item 时间偏移事件的中位偏移量 $\sim\pm$15~ms（毫秒级，非 wrap 级错误）；
  \item 1B 比 1K 在每个恢复区域多出约 100--2800 个事件，
    说明 1K 在 FIFO 复位边缘做了更激进的事件过滤；
  \item 几乎没有仅存在于 1K 中的事件（0--37 个）。
\end{itemize}

结论：1B 时间重建保留了更完整的事件集，
剩余差异来自 1K 管线的过滤策略而非 1B 的时间错误。

\subsubsection{FIFO 复位后向修正效果}

以 Box~B 为例，后向 SEC 修正显著改善了 FIFO 复位边界处的事件归位：
\begin{itemize}
  \item $T+249$~s：修正前 402/9960 ($-96\%$) $\rightarrow$ 修正后 10065/9960 ($+1.1\%$)；
  \item $T+277/278$~s：双重模糊区域（两个锚点的 wrap 边界重合），
    约 3000 个事件仍有 $\sim P_\mathrm{wrap}$ 偏移。
\end{itemize}

\subsubsection{光变曲线重建}

%% Figure 6: 光变曲线重建效果
\begin{figure*}[t]
\figplaceholder{图~6：GRB~221009A 光变曲线重建效果对比。
左列：原始观测光变曲线（三个 Box）。
右列：重建后光变曲线（FIFO reset gap 补全 + 静默丢数补全 + 深度饱和修正）。
红色：observed；蓝色：filled\_gap（FIFO reset 补全）；
绿色：filled\_sd（静默丢数 + 深度饱和修正）。}
\caption{GRB~221009A（极端饱和）的光变曲线重建效果。}
\label{fig:reconstruction}
\end{figure*}

GRB~221009A 检测到 138 个静默丢数候选（Box~A 42、Box~B 42、Box~C 39），
分布在约 90 个独立 CCSDS 包中。
静默丢数位置集中在包头和包尾（事件间隔索引 0 和 106），
对应 MCU 在包边界开始/结束读取 FIFO 的时刻。

三机箱交叉参考的覆盖率约 92\%（至少一个参考 box 可用），
仅约 8\% 的 gap 三个机箱同时饱和，此时退化为均匀分配。

%% Figure 7: 三机箱覆盖率统计
\begin{figure}[t]
\figplaceholder{图~7：GRB~221009A 三机箱交叉参考覆盖率统计。
柱状图：每个 Box 的 FIFO reset gap 中，有 0/1/2 个参考 box 可用的比例。}
\caption{三机箱交叉参考覆盖率统计。}
\label{fig:coverage}
\end{figure}

\subsection{1B/1K 对比统计}\label{sec:compare}

表~\ref{tab:grbs} 汇总了三个 GRB 的 1B/1K 对比。
轻度饱和事件（200415A、260226A）的 1B/1K 完美一致，
验证了时间重建算法的正确性。
极端饱和事件（221009A）中 1B 系统性多出少量事件，
这是因为 1B 保留了更完整的事件集（1K 在 FIFO 复位边缘有额外过滤）。

Box~C 的 3.5\% 差异主要集中在 $T+510$~s 附近（1B/1K 比值 $= 2.04$），
经分析系 1K 管线在此处丢失数据（1B 比 1K 多约 12 万事件），
属于 1K 管线的问题而非 1B 重建的误差。

%% ======================================================================
\section{讨论}\label{sec:discussion}

\subsection{4-bit CRC 碰撞}

1B 数据的事例校验仅使用 4-bit CRC，碰撞概率约 $1/16 \approx 6.3\%$。
在高计数率下，随机字节序列有一定概率通过 CRC 校验产生"幽灵事件"。
这些幽灵事件的 ptime 和 channel 值是随机的，
会干扰包内中位数 ptime 的计算，影响 wrap tracking 的判定。

当前的保护措施是设置 \texttt{MIN\_EVENTS\_FOR\_MEDIAN} $= 50$：
有效事件数不足 50 的包被视为损坏包，跳过 wrap 检测和中位数更新。
在 GRB~221009A 的约 32 万个包中仅发现 2 个严重损坏包（0.001\%），
说明传输链路整体可靠。

\subsection{双重模糊}

当 FIFO 复位后过期锚点和新鲜锚点的 anchor ptime 差异恰好使同一批事件
在两个方向上都处于 wrap 边界时，前后对照法无法消歧。
表现为两个候选的残差 $|e_a - e_b| \approx P_\mathrm{wrap}$。

Box~B $T+277/278$~s 附近约 3000 个事件即属此情况。
从数据本身无法进一步消歧，需要外部约束（如 1K 管线对照）辅助判断。
此类情况发生的概率取决于过期锚点和新鲜锚点 ptime 的偶然关系，
总体影响有限。

\subsection{三机箱同时饱和}

约 8\% 的 FIFO reset gap 三个机箱同时饱和，此时交叉参考不可用，
光变曲线重建退化为均匀分配（假设 gap 内计数率恒定）。
这主要发生在 GRB 光变曲线峰值附近、所有机箱同时经历极高计数率的短暂时段。

\subsection{包间中等 gap}

当前 $F_\mathrm{gap} = 100$ 的阈值下，包间 gap 为正常值 10--50 倍
但不足 100 倍的区间不会被检测为 FIFO 复位。
这些可能是未触发完全复位的 FIFO 部分拥塞事件。
降低 $F_\mathrm{gap}$ 可以捕捉更多，但会增加误报率。
在当前实现中，这些中等 gap 可能在光变曲线上表现为窄缝。

\subsection{适用范围}

本文方法适用于所有 HXMT HE 的 1B 级数据，
但饱和分析和光变曲线重建的效果取决于饱和程度：
\begin{itemize}
  \item 非饱和或轻度饱和：时间重建高精度，光变曲线几乎无损；
  \item 中度饱和：FIFO reset gap 可通过交叉参考较好地重建；
  \item 极端饱和（如 221009A 峰值）：深度饱和区 $R_\mathrm{true}$ 估算
    受限于 MCU 读模式，三机箱同时饱和时重建为均匀分配。
\end{itemize}

%% ======================================================================
\section{结论}\label{sec:conclusion}

本文发展了一套完整的 HXMT HE 1B 级数据饱和分析与光变曲线重建方法，
主要结论如下：

\begin{enumerate}
  \item 四遍时间重建算法通过 wrap tracking、SEC 锚点和后向修正等机制，
    有效克服了 FIFO 拥塞和复位导致的时间戳失真问题。
    在非饱和区与 1K 管线结果完全一致（$\Delta = 0$），
    在极端饱和区达到 97.9\%--100\% 的逐事件精确匹配率。

  \item 识别了 HE 饱和丢数的三种模式
    ——包内无痕丢数、FIFO 复位整包丢失和深度饱和——
    并分别发展了对应的检测与重建方法。

  \item 三机箱交叉参考策略利用独立机箱间的互补性，
    在 92\% 的 FIFO reset gap 中提供了至少一个可靠的参考光变曲线形状。

  \item 1B 级处理保留了比 1K 管线更完整的事件集，
    1K 管线在 FIFO 复位边界的激进过滤导致了约 0.003\%--3.5\% 的事件缺失，
    而 1B 重建避免了这些额外损失。

  \item 已知局限包括 4-bit CRC 碰撞产生的幽灵事件（$\sim$6\%）、
    双重模糊区域的 wrap 判定残余（影响 $\sim$3000 事件/$\sim$400 万事件）、
    以及三机箱同时饱和时交叉参考的不可用（$\sim$8\% 的 gap）。
\end{enumerate}

该方法以 Rust 语言实现为命令行工具，可处理任意 HXMT HE 1B 数据文件，
为利用 1B 数据开展饱和区域的高时间分辨率科学分析提供了可靠工具。

\begin{acknowledgments}
感谢慧眼卫星团队提供数据支持。
本研究使用了中国科学院高能物理研究所的计算资源。
\end{acknowledgments}

\end{document}
